\chapter{Discussion}
\label{ch:discussion}

\section{What the Results Mean}

\paragraph{The task is learnable, and not yet solved by general models.} A fine-tuned
convolutional network names the statistical finding of a real chart correctly almost
nine times in ten, while the strongest open vision--language model we tested manages
fewer than one in two without task training. The gap is not about recognising the chart
type (a rule that uses only chart type reaches 0.288): it is about judging shape, spread
and change, which is exactly what the verified labels supervise. This supports the
central premise of the thesis (Chapter~\ref{ch:intro}): a dataset whose labels state
what the data \emph{shows}, rather than what a question asks, trains a capability that
question-answering corpora do not.

\paragraph{Most remaining errors are judgements of degree.} The confusions that remain
(bimodal against right-skewed, heavy-tailed against right-skewed, significant against
no group difference) are between classes separated by a threshold on a continuous
statistic. Near a threshold, the chart itself is ambiguous: a right-skewed histogram
with a shoulder can look bimodal to a person too. Some of these errors are therefore
likely irreducible from pixels at this resolution; the oracle reaches 0.9995 only
because it reads the numbers.

\paragraph{Where the training data came from mattered as much as the model.} The three
largest improvements in this thesis all came from the data rather than from
architecture: removing the rendering leak (which made the problem honest), adding real
training charts, and drawing training charts with more than one library. Architecture
changes within the CNN family moved real accuracy by one to two points.

\section{Lessons from the Rendering Leak}

The leak is the most important methodological lesson of the project, and we report it
in full because it could affect other synthetic chart datasets.

\begin{enumerate}[leftmargin=1.4em]
  \item \textbf{A helpful annotation can be a label.} The trend line, change-point line
  and anomaly circles were added to make charts more informative, and they were drawn
  only where the finding held. Any rendering step that looks at the label is a potential
  leak, however natural it seems.
  \item \textbf{A data audit does not replace a model audit.} Eleven shortcut tests on
  metadata and pixel statistics all passed while trained models were reading the marks.
  The failure was only visible by changing the input a trained model sees and measuring
  what happened.
  \item \textbf{A class that cannot carry the shortcut is a control.} Seasonal series
  never carried a mark, and their accuracy did not move when marks were removed. That
  single class separated ``the models read the marks'' from ``the renders changed''.
  \item \textbf{High accuracy on an easy-looking class is a warning.} Level-shift charts
  scored 1.000 with marks and 0.355 without. A perfect score on a class that should be
  hard deserves suspicion, not celebration.
\end{enumerate}

\section{Threats to Validity}

\paragraph{Internal validity.} All selection used the validation split, but the
validation split has only 553 real charts, so choices made on it are noisy; this is why
the deployed ensemble uses equal weights and why alternatives are reported. The choice to
deploy the full-recipe models was made after all models were trained; the alternatives
(averaging every model, 0.862; greedy selection, 0.865) bound its effect at under a
point. Kaggle probability outputs were checked against a local re-computation (100\%
prediction agreement) before use.

\paragraph{Construct validity.} The labels encode conventional thresholds
($|r|>0.7$ for ``strong'', $\varepsilon^2>0.14$ for ``large''). A model that scores well
has learned these conventions, which are reasonable but not universal. The evidence
boxes used for grounding mark where the evidence is by construction, not where a human
would look.

\paragraph{External validity.} Real test charts come from six public repositories and
are drawn by our renderers, not taken from papers or reports in the wild. The Vega-Lite
test set measures an unseen renderer, not unseen chart conventions (hand-drawn charts,
screenshots, charts with annotations of their own). Real class frequencies are uneven,
and four classes have fewer than 20 real test charts, so their accuracies have wide
uncertainty.

\paragraph{Statistical conclusion validity.} Every headline accuracy carries a bootstrap
confidence interval, and system comparisons use paired exact McNemar tests on the same
charts. Per-class accuracies for rare classes are reported with their $n$ and should not
be over-read. Per-class results for the rarest real classes are indicative only: with $n=13$ to $26$ real test charts, a class's 95\% interval spans roughly $\pm 19$ to $\pm 27$ points. The confidence intervals reflect test-set sampling; training variance across random seeds is
reported separately (Section~\ref{sec:rq1}).

\section{Limitations}

\begin{enumerate}[leftmargin=1.4em]
  \item \textbf{Mostly synthetic.} 69.6\% of the balanced corpus is synthetic; only ten
  clean inverted-U relationships exist in 719 real tables.
  \item \textbf{Single label.} Charts showing two findings were excluded, so models will
  not report a trend and an anomaly together.
  \item \textbf{A small generative model.} The generative result comes from a 2B-parameter
  model trained for one epoch on 4,000 samples and scored on 300 charts (82 real), so its
  real-chart accuracy has a wide interval.
  \item \textbf{Small human study.} Two annotators, one an author, checked 188 labels
  (Section~\ref{sec:human}); they confirmed 73.7\%, but only 25--40\% for four classes whose
  findings are statistically true yet hard to see. No human classification baseline was
  completed.
  \item \textbf{Compute.} All training used free T4 GPUs under a weekly quota, which
  limited the number of models, epochs and seeds; results come from single training runs
  per configuration, without seed variance.
\end{enumerate}
